\subsection{Partial spectra and sparse computation}\label{sec:sparse}
The preceding studies use full dense spectra. A separate extension stores
the input matrices in compressed sparse row format and computes a partial
spectrum through ARPACK, using SciPy's \texttt{eigsh} interface
\cite{scipyeigsh}. The initial request is $\min(n,\max(r+8,2r))$ vectors;
it expands when the returned boundary is too close to the rank cutoff.
QR and Rayleigh--Ritz refinement, residual checks, and a 32 MiB spectrum
cache accompany the same reusable endpoint search. Dense fallbacks are
explicitly recorded and disabled above $n=1024$. This is an implementation
of established sparse spectral methods, not a new eigensolver.

The comparison uses a dense \emph{partial}-spectrum LAPACK backend with
the same outer search and recovery. Both receive identical synthetic
positive-semidefinite matrices: independently weighted graph Laplacians
plus positive diagonal terms. Edge weights lie in $[0.05,0.25]$, reward
diagonals in $[0.2,2]$, resource diagonals in $[0.1,1.5]$, and overheads in
$[0,0.15]$. The graph is a row-major grid of width $\lfloor\sqrt n\rfloor$,
possibly with an incomplete final row. These are computational test
families without a calibrated physical interpretation.

The initial sparse study stopped after 18 of 72 planned runs: both backends
failed all three repetitions of one rank-two case at $n=128$. A narrow
avoided eigenvalue crossing caused price bisection to reach binary64
resolution before meeting the $10^{-9}$ endpoint tolerance. The revised
implementation rotates the adjacent cutoff vectors to match the required
chord slope, evaluates the original primal objective, and accepts only
under the unchanged gap or decision criterion. The six failures and their
source snapshot are retained. This recovery is applied to both backends;
the failed study's seeds are thereafter treated as development data.

A new protocol uses fresh seeds 6400 and 6401, $n=128,256,512$,
$r=2,8$, two sites, and $K=2$. Three repetitions and two backends give
72 timed runs in 12 paired cells, with alternating method order, one BLAS
thread, and the preceding tolerances. Every saved menu is checked with
independent full dense eigenvalue evaluation at all affine-line crossings
and covering-segment endpoints. Those numerical cross-checks are distinct
from outward certificates. Table~\ref{tab:sparse} reports per-cell median
times; smaller instances need not benefit from sparse computation.

\begin{table}[t]\centering\small
\caption{Sparse confirmation. Times are medians of three runs in seconds;
the ratio is dense partial divided by CSR/ARPACK. A ratio above one favors
the sparse backend. Each row has two sites and two presets.}
\label{tab:sparse}\begin{tabular}{rrrrrr}\toprule
$n$ & $r$ & Seed & Dense partial & CSR/ARPACK & Ratio\\\midrule
128 & 2 & 6400 & 1.60 & 6.63 & 0.24 \\
128 & 2 & 6401 & 1.24 & 5.01 & 0.25 \\
128 & 8 & 6400 & 2.08 & 7.67 & 0.27 \\
128 & 8 & 6401 & 2.01 & 7.16 & 0.28 \\
256 & 2 & 6400 & 4.97 & 9.72 & 0.51 \\
256 & 2 & 6401 & 4.24 & 6.99 & 0.61 \\
256 & 8 & 6400 & 6.22 & 10.71 & 0.58 \\
256 & 8 & 6401 & 6.11 & 11.16 & 0.55 \\
512 & 2 & 6400 & 23.64 & 12.43 & 1.90 \\
512 & 2 & 6401 & 22.08 & 11.75 & 1.88 \\
512 & 8 & 6400 & 26.02 & 14.25 & 1.83 \\
512 & 8 & 6401 & 40.09 & 28.81 & 1.39 \\
\bottomrule\end{tabular}

\end{table}

All 72 confirmation runs completed, and all 24 saved menus passed the
independent numerical check. The median dense-partial/CSR runtime ratio
across the 12 cells is $0.564$. Its per-dimension medians
are $0.259$, $0.564$, and
$1.853$ for $n=128,256,512$, respectively.
Thus the small-instance overhead must be weighed against the measured
larger-instance behavior, rather than reporting only a favorable size.
The maximum difference between reported and independently evaluated losses
over confirmation and scaling menus is $2.5\times10^{-14}$.
The sparse runs in these two studies used 0 dense fallbacks.

Table~\ref{tab:sparse-scaling} extends the numerical screen to $n=4096$.
For the single $n=1024$ grid, the dense-partial/CSR ratio is $7.59$;
this one observation is not a repeated speedup estimate. Independent
ordered tridiagonal, full dense, and banded LAPACK calculations check path,
$n=1024$ grid, and $n=4096$ grid menus, respectively. The banded check uses
the known grid bandwidth, not the ARPACK proposals.

Outward verification and standalone replay also succeed for the four
$n=128$ CSR grid menus and all four path menus in the scaling screen,
each with a global bound width $10^{-5}$ to displayed precision. The grid certificates
use dense arithmetic; the path certificates use the tridiagonal verifier.
The banded verification study below extends certification to the larger grids.
Recorded array storage includes all cached preset bases, which can exceed
the spectrum-cache budget; it excludes transient solver workspaces and
Python overhead and is not a peak-memory measurement.


Partial eigenpair residuals do not establish extremal ordering or complete
cutoff multiplicity. The grid results therefore use separate ordered
spectral checks, with outward dense verification on selected tractable
cases. For paths, an outward Sturm-count verifier certifies ordered
eigenvalue sums without a dense matrix, using classical tridiagonal
inertia\cite{demmel1993}. Appendix~\ref{app:sturm} gives the argument.
The descriptive larger-instance screen also compares paths against a
specialized tridiagonal LAPACK solver; general CSR performance should not
be inferred from beating an unnecessarily dense path implementation.

\begin{table}[t]\centering\small
\caption{Descriptive scaling screen: one solve per method and case,
$r=2$, two sites, $K=2$. These are single observations, not repeated
confirmation estimates. Retained memory counts input, spectrum-cache peak,
and final basis-sample arrays, not peak process RSS.}
\label{tab:sparse-scaling}\begin{tabular}{lrlrr}\toprule
Graph & $n$ & Backend & Seconds & Retained MiB\\\midrule
Path & 1024 & CSR/ARPACK & 17.09 & 60.6 \\
Path & 1024 & Tridiagonal & 5.03 & 54.5 \\
Path & 4096 & CSR/ARPACK & 64.34 & 140.8 \\
Path & 4096 & Tridiagonal & 23.98 & 140.8 \\
Grid & 1024 & Dense partial & 134.49 & 83.6 \\
Grid & 1024 & CSR/ARPACK & 17.73 & 59.8 \\
Grid & 4096 & CSR/ARPACK & 189.80 & 356.3 \\
\bottomrule\end{tabular}

\end{table}

\paragraph{Outward verification of every CSR grid menu.}
A further verification study certifies all 12 CSR confirmation menus and
both CSR grid scaling menus, through $n=4096$, using interval $LDL^T$
elimination in band storage. All 14 certificates pass standalone replay,
with global gaps below $1.01\times10^{-5}$ (Table~\ref{tab:banded}).
This strengthens the preceding numerical checks without changing any
optimizer, saved menu, or isolated runtime measurement.

An initial binary64 interval implementation rejected a rank-eight witness
because its eigenvalue brackets became too wide to prove exclusion. That
attempt is retained. The revised verifier retries unresolved counts at
$n\le512$ with directed 50-digit decimal intervals; 57 counts use this
path. Both larger grid cases need only outward binary64 arithmetic.
The final checker independently encloses ordered eigenvalue sums and actual
polar-orthonormalized projectors; it imports no menu solver. Six CPU
processes check independent price/site pairs. These parallel verification
times are separate from the isolated optimization study.

For half-bandwidth $b$, each unpivoted inertia count uses $O(nb^2)$ work
and $O(nb+b^2)$ storage. The grid has $b=\lfloor\sqrt n\rfloor$;
at $n=4096$, packed input bands occupy about 6.1 MiB and an interval factor
about 4.1 MiB per worker, excluding temporary arrays and process overhead.
This is an advantage of band structure, not a guarantee for arbitrary
sparse graphs. Ambiguous pivots cause rejection; high bandwidth can erase
the storage benefit. Appendix~\ref{app:sturm} gives the factorization
argument and arithmetic assumptions.

\begin{table}[t]\centering\small
\caption{Independently replayed outward certificates for CSR grid menus.
Decimal counts record successful higher-precision inertia retries;
the displayed gap column is rounded. All unrounded gaps are below
$1.01\times10^{-5}$.}
\label{tab:banded}\begin{tabular}{rrrrr}\toprule
$n$ & Menus & Half-bandwidth & Decimal counts & Max. gap / $10^{-5}$\\\midrule
128 & 4 & 11 & 18 & 1.000029 \\
256 & 4 & 16 & 26 & 1.000031 \\
512 & 4 & 22 & 13 & 1.000031 \\
1024 & 1 & 32 & 0 & 1.000008 \\
4096 & 1 & 64 & 0 & 1.000008 \\
\bottomrule\end{tabular}

\end{table}

